% 20-stdlib-impl.tex — Chapter 20: Standard Library Implementation
\chapter{Standard Library Implementation}
\label{chap:20-stdlib-impl}
\index{standard library}

\pnum
The Lain standard library consists of a set of \texttt{.ln} source files under
\texttt{std/}.  These files are written in Lain itself and are imported by user
programs via \lain{import std.io}, \lain{import std.math}, etc.  The standard
library is compiled along with user code in the same single-pass pipeline; there
is no pre-compiled library archive.

% -------------------------------------------------------------------------
\section{std.c --- C interoperability}
\label{sec:stdlib-c}
\index{std.c}
% -------------------------------------------------------------------------

\pnum
\srcref{std/c.ln} declares the C standard library interface using
\lain{extern proc} and \lain{extern func} declarations.  These declarations
make C functions available to Lain code without requiring a \lain{c\_include}
directive in user code.  The module also declares the opaque \texttt{FILE} type
via \lain{extern type FILE}.

% -------------------------------------------------------------------------
\section{std.io --- Console I/O}
\label{sec:stdlib-io}
\index{std.io}
% -------------------------------------------------------------------------

\pnum
\srcref{std/io.ln} provides basic console output.  It wraps the C
\texttt{printf} and \texttt{puts} functions:

\begin{center}
\begin{tabular}{ll}
\toprule
\textbf{Function} & \textbf{Description} \\
\midrule
\lain{print(s u8[:0])}   & Print a string without newline. \\
\lain{println(s u8[:0])} & Print a string followed by a newline. \\
\bottomrule
\end{tabular}
\end{center}

% -------------------------------------------------------------------------
\section{std.fs --- File system}
\label{sec:stdlib-fs}
\index{std.fs}
% -------------------------------------------------------------------------

\pnum
\srcref{std/fs.ln} provides basic file I/O.  The \lain{File} type is linear
(owned): it must be explicitly closed after use.

\begin{center}
\begin{tabular}{lll}
\toprule
\textbf{Function} & \textbf{Signature} & \textbf{Description} \\
\midrule
\lain{open}  & \lain{proc(path u8[:0]) mov File} & Open file for writing. \\
\lain{close} & \lain{proc(f mov File)}            & Close and consume the file handle. \\
\lain{write} & \lain{proc(f mut File, s u8[:0])}  & Write a string to a file. \\
\bottomrule
\end{tabular}
\end{center}

\pnum
The \lain{File} type is declared as \lain{extern type} backed by a C \texttt{FILE *}.
Opening returns a \lain{mov File} (owned), ensuring the caller cannot forget to
close it without a compiler error.

% -------------------------------------------------------------------------
\section{std.math --- Integer math}
\label{sec:stdlib-math}
\index{std.math}
% -------------------------------------------------------------------------

\pnum
\srcref{std/math.ln} provides pure integer math utilities implemented directly
in Lain:

\begin{center}
\begin{tabular}{ll}
\toprule
\textbf{Function} & \textbf{Description} \\
\midrule
\lain{min(a int, b int) int}                & Minimum of two integers. \\
\lain{max(a int, b int) int}                & Maximum of two integers. \\
\lain{abs(x int) int}                       & Absolute value. \\
\lain{clamp(x int, lo int, hi int) int}     & Clamp \texttt{x} to $[\textit{lo}, \textit{hi}]$. \\
\bottomrule
\end{tabular}
\end{center}

\pnum
All four functions are declared \lain{func} (pure): they contain no loops and
return immediately, satisfying the termination guarantee.

% -------------------------------------------------------------------------
\section{std.string --- String utilities}
\label{sec:stdlib-string}
\index{std.string}
% -------------------------------------------------------------------------

\pnum
\srcref{std/string.ln} provides string operations over \lain{u8[:0]} slices
(NUL-sentinel strings).  Functions include length computation, substring
extraction, and comparison.

% -------------------------------------------------------------------------
\section{std.option and std.result}
\label{sec:stdlib-option-result}
\index{std.option}
\index{std.result}
% -------------------------------------------------------------------------

\pnum
\srcref{std/option.ln} and \srcref{std/result.ln} define generic algebraic
data types using comptime generics:

\begin{laincode}
enum Option(comptime T type) {
    Some { value T }
    None
}

enum Result(comptime T type, comptime E type) {
    Ok  { value T }
    Err { error E }
}
\end{laincode}

\pnum
These are the standard vocabulary types for optional and fallible computations.
Instantiation is monomorphic: each distinct type argument combination produces a
separate C struct in the output.
